% !PRINTED 653
\weberpage{653}
%<<CONT
the function ``1'' must also be representable. Let us assume that
$\omega_1$ is not divisible by $\pi$. By \S\ 182, 6. there are constants
$c_1, c_2, \ldots, c_n$, the first of which may, without prejudice to generality, be
assumed $= 1$, and which satisfy the conditions:
\[
\omega_1 \equiv 1, \ \omega_2 \equiv c_2, \ \ldots, \ \omega_n = c_n \quad (\mathrm{mod}\ \pi)
% [sic] printed $\omega_n = c_n$ with = instead of \equiv, the relation is a congruence modulo $\pi$ like the two preceding ones
\]
We set:
\weq{14}{
\begin{aligned}
\alpha_1 &= (z - c)\,\omega_1,\\
\alpha_2 &= \omega_2 - c_2\,\omega_1,\\
&\cdot\ \cdot\ \cdot\ \cdot\ \cdot\ \cdot\ \cdot\\
\alpha_n &= \omega_n - c_n\,\omega_1.
\end{aligned}
}

That this is a basis of $\pi$ one sees at once if one represents
the function (13) in this way:
\weq{15}{
\omega = x_2\alpha_2 + \cdots + x_n\alpha_n + (x_1 + c_2x_2 + \cdots + c_nx_n)\,\omega_1\,;
}
since $\alpha_2, \ldots, \alpha_n$ are divisible by $\pi$, $\omega$ can be divisible by
$\pi$ only if the integral function in $Z$:
\[
x_1 + c_2x_2 + \cdots + c_nx_n
\]
is divisible by $(z - c)$, and consequently $\omega$ is representable by formula (5).
Hence formulae (10) and (14) give:
\weq{16}{
N_a(\pi) = z - c\,,
}
and thus the theorem:

2. \gesperrt{The absolute norm of a prime functional is a linear function of $z$.}

This theorem marks an essential difference between
the theories of numbers and of functions. In number theory
 there are prime ideals of various degrees whose norm is a
power of a prime. In the theory of functions we have
only prime functionals of the first degree. This comes about because
by virtue of the fundamental theorem of algebra we can decompose every integral
function of one variable, with respect to the modulus $\pi$, into linear
factors.

The number of prime factors into which $(z - c)$ splits is accordingly
 $n$, since $N(z-c) = (z-c)^n$.

If we collect the mutually equal prime factors into
powers, then
\weq{17}{
z - c = \varepsilon\,\pi_1^{e_1}\pi_2^{e_2}\pi_3^{e_3}\ldots\,,
}
in which
\weq{18}{
n = e_1 + e_2 + e_3 + \cdots\,,
}
where $\varepsilon$ is a unit.
